\chapter{Tests and non-vacuity}\label{chap:tests}

\section*{At a glance}
\begin{itemize}
\item \lead{Goal} Show that the hypotheses of the main statements can be discharged on concrete
  programs, relate the restated relations to lean4lean's, and pin the design choices of the
  divergence register by tests.
\item \lead{Main results} The final theorem \code{erase\_correct} on NV-1, through the
  kernel-checked run of \code{\#erase}; seven instances of \code{erases\_correct} (NV-1 to NV-7),
  each with a unique non-$\Box$ witness, NV-3 and NV-4 through recursive constants; the approved
  statement of \code{erase\_correct}, which check C12 compares with the theorem; \code{TrS} and
  \code{ProgEnv} are instances of lean4lean's \code{TrExprS} and \code{TrEnv'}; the register
  tests of DV-11 and DV-18; two instances of \code{Pure.isErasable\_sound}.
\item \lead{Status} \stProved{} every node; \stInherited{} the instances of \code{erases\_correct}
  and \code{erase\_correct} and the approved statement inherit L1--L6 from their theorem, those
  of oracle soundness L1--L5, and the bridge tests \code{TrProj}.
\item \lead{Nodes} \bpnodes{chap:tests}
\item \lead{Lean files}
  \code{proof/EraseProof/Test/\{NV1, NV1Run, NV2, NV3, NV4, NV5, NV6, NV7, Stub, Bridge, LBEval,
  Atoms, Oracle\}.lean}; namespace \code{EraseProof.Test}.
\end{itemize}

\section{Non-vacuity instance NV-1}\label{sec:tests-nv1}

\code{erases\_correct} on \code{one (fun a : A => a)}: the source evaluation takes a $\delta$, a
$\beta$ and three atom steps. Section~\ref{sec:tests-nv1-run} runs \code{\#erase} on it.

\begin{definition}[The program of NV-1]
  \label{def:EraseProof-Test-NV1-decls0}
  \lean{EraseProof.Test.NV1.decls0, EraseProof.Test.NV1.e0, EraseProof.Test.NV1.tyA,
    EraseProof.Test.NV1.A_val, EraseProof.Test.NV1.cnE, EraseProof.Test.NV1.oneE,
    EraseProof.Test.NV1.CN_val, EraseProof.Test.NV1.one_val, EraseProof.Test.NV1.venv0,
    EraseProof.Test.NV1.e0', EraseProof.Test.NV1.vA, EraseProof.Test.NV1.ty1,
    EraseProof.Test.NV1.vAA, EraseProof.Test.NV1.cnV, EraseProof.Test.NV1.oneV,
    EraseProof.Test.NV1.Aconst, EraseProof.Test.NV1.CNdef, EraseProof.Test.NV1.onedef,
    EraseProof.Test.NV1.envA, EraseProof.Test.NV1.envCN0, EraseProof.Test.NV1.envCN,
    EraseProof.Test.NV1.envOne0, EraseProof.Test.NV1.view0, EraseProof.Test.NV1.σ0,
    EraseProof.Test.NV1.v0}
  \leanok
  \uses{def:EraseProof-evalEnvOf}
  \srcloc{proof/EraseProof/Test/NV1.lean}{57}
  \lead{In short} \code{axiom A : Type}, \code{def CN := (A -> A) -> A -> A}, \code{def one : CN :=
  fun s z => s z}, and the erased term \hl{\code{one (fun a : A => a)}}.
  \begin{itemize}
  \item \code{decls0}, \code{e0}: the declarations, newest first, and the term;
  \item \code{venv0}, \code{e0'}: their lean4lean model and translation, built step by step
        (\code{envA}, \code{envCN0}, \code{envCN}, \code{envOne0});
  \item \code{view0}, \code{sigma0}: a list view and its evaluation environment; \code{v0}: the
        value \code{fun z => (fun a => a) z}.
  \end{itemize}
\end{definition}

\begin{definition}[The erased program of NV-1]
  \label{def:EraseProof-Test-NV1-lenv0}
  \lean{EraseProof.Test.NV1.lenv0, EraseProof.Test.NV1.oneL, EraseProof.Test.NV1.t0,
    EraseProof.Test.NV1.v0'}
  \leanok
  \uses{def:GlobalDecl, def:LBTerm}
  \srcloc{proof/EraseProof/Test/NV1.lean}{247}
  \lead{In short} The $\lambda_\Box$ environment \code{lenv0} holds \hl{\code{one} with its erased
  body}, as \code{\#erase} emits it.
  \begin{itemize}
  \item \code{t0}: \code{one (fun a => a)}, the erasure of \code{e0};
  \item \code{v0'}: \code{fun z => (fun a => a) z}, the erasure of \code{v0}.
  \end{itemize}
\end{definition}

\begin{lemma}[NV-1: the model-side hypotheses]
  \label{lem:EraseProof-Test-NV1-henv}
  \lean{EraseProof.Test.NV1.henv, EraseProof.Test.NV1.he, EraseProof.Test.NV1.venv0_A,
    EraseProof.Test.NV1.hAty, EraseProof.Test.NV1.hAAty, EraseProof.Test.NV1.cnV_ty,
    EraseProof.Test.NV1.cn_unfold, EraseProof.Test.NV1.trCN, EraseProof.Test.NV1.trOne,
    EraseProof.Test.NV1.oneWF, EraseProof.Test.NV1.pCN}
  \leanok
  \uses{def:EraseProof-Test-NV1-decls0, def:EraseProof-ProgEnv, def:EraseProof-TrS}
  \srcloc{proof/EraseProof/Test/NV1.lean}{173}
  \lead{In short} \code{henv : ProgEnv decls0 venv0} and \code{he : TrS venv0 [] [] e0 e0'}:
  \hl{the environment and typing hypotheses} of the main statements hold on NV-1.
\end{lemma}
\begin{proof}
  \uses{def:EraseProof-ProgEnv}
  \leanok
  \lead{Idea} One \code{ProgEnv} rule per declaration; typing derivations in lean4lean's
  \code{HasType}, with \code{CN} unfolded by its defining equation.
\end{proof}

\begin{lemma}[NV-1: the view agrees with the program]
  \label{lem:EraseProof-Test-NV1-hview}
  \lean{EraseProof.Test.NV1.hview}
  \leanok
  \uses{def:EraseProof-Test-NV1-decls0, def:EraseProof-ViewAgrees}
  \srcloc{proof/EraseProof/Test/NV1.lean}{198}
  \lead{In short} \code{hview : ViewAgrees view0 decls0}: \hl{the view hypothesis of
  \code{erase\_correct}} holds on NV-1.
\end{lemma}
\begin{proof}
  \uses{def:EraseProof-ViewAgrees}
  \leanok
  \lead{Idea} One case per declaration; each lookup is by \code{rfl}.
\end{proof}

\begin{lemma}[NV-1: the relation-level instance]
  \label{lem:EraseProof-Test-NV1-concl}
  \lean{EraseProof.Test.NV1.concl, EraseProof.Test.NV1.inst, EraseProof.Test.NV1.witness,
    EraseProof.Test.NV1.hev,
    EraseProof.Test.NV1.hsub, EraseProof.Test.NV1.mem_names, EraseProof.Test.NV1.hinj,
    EraseProof.Test.NV1.lookup0, EraseProof.Test.NV1.hlc, EraseProof.Test.NV1.hblocks,
    EraseProof.Test.NV1.trAA, EraseProof.Test.NV1.erOne, EraseProof.Test.NV1.her,
    EraseProof.Test.NV1.hdeps, EraseProof.Test.NV1.lbev, EraseProof.Test.NV1.herv}
  \leanok
  \uses{thm:EraseProof-erases_correct, lem:EraseProof-Test-NV1-henv,
    def:EraseProof-Test-NV1-lenv0}
  \srcloc{proof/EraseProof/Test/NV1.lean}{319}
  \inherited{L1, L2, L3, L4, L5, L6}
  \lead{In short} \code{inst} applies \code{erases\_correct} on NV-1 with \hl{every hypothesis a
  checked term}; \code{concl} is its conclusion at the $\lambda$ \code{v0'}.
  \begin{itemize}
  \item \code{hev}, \code{hsub}, \code{hinj}, \code{hlc}, \code{hblocks}, \code{her},
        \code{hdeps}: the hypotheses besides \code{henv} and \code{he};
  \item \code{lbev}, \code{herv}: \code{t0} evaluates to \code{v0'}, which erases \code{v0};
  \item \code{witness}: every witness of the conclusion is \code{v0'}.
  \end{itemize}
\end{lemma}
\begin{proof}
  \uses{thm:EraseProof-erases_correct, lem:EraseProof-Test-LBEval-deterministic}
  \leanok
  \lead{Idea} Apply the theorem; \code{witness} by determinism of $\lambda_\Box$ evaluation; the
  footprint is the theorem's.
\end{proof}

\section{NV-1 through \texorpdfstring{\code{\#erase}}{\#erase}}\label{sec:tests-nv1-run}

\code{erase\_correct} on NV-1: the kernel checks the run of \code{\#erase}'s entry point on
\code{e0}, one traversal call at a time, and the final theorem applies with every hypothesis a
checked term.

\begin{definition}[The run data of NV-1]
  \label{def:EraseProof-Test-NV1-p0}
  \lean{EraseProof.Test.NV1.p0, EraseProof.Test.NV1.decls0', EraseProof.Test.NV1.stRun,
    EraseProof.Test.NV1.lS, EraseProof.Test.NV1.lZ, EraseProof.Test.NV1.lA}
  \leanok
  \uses{def:EraseProof-Test-NV1-decls0, def:EraseProof-Test-NV1-lenv0, def:ASTType,
    def:Erasure-EraseT, def:Erasure-TravCtx, def:EraseProof-ExprSpec}
  \srcloc{proof/EraseProof/Test/NV1Run.lean}{320}
  \lead{In short} \code{p0 = untyped lenv0 (some t0)}: \hl{the program \code{\#erase} emits for
  NV-1}.
  \begin{itemize}
  \item \code{decls0'}: the closure of \code{e0}, in \code{collectDeps}' order \code{A},
        \code{CN}, \code{one};
  \item \code{stRun}: the traversal state after the run, with \code{one} registered and the
        $\lambda_\Box$ environment \code{lenv0};
  \item \code{lS}, \code{lZ}, \code{lA}: the locals of the binders \code{s}, \code{z}, \code{a},
        at the counter's first three variables.
  \end{itemize}
\end{definition}

\begin{lemma}[NV-1: the closure and the route]
  \label{lem:EraseProof-Test-NV1-hcollect}
  \lean{EraseProof.Test.NV1.hcollect, EraseProof.Test.NV1.eraseEntry_of_erasePure}
  \leanok
  \uses{def:Erasure-collectDeps, def:Erasure-eraseEntry, def:Erasure-erasePure,
    def:EraseProof-Test-NV1-decls0, def:EraseProof-Test-NV1-p0}
  \srcloc{proof/EraseProof/Test/NV1Run.lean}{28}
  \lead{In short} \code{collectDeps view0 e0 = ok decls0'}, so \hl{\code{eraseEntry} on NV-1
  returns what \code{erasePure} returns over \code{decls0'}}.
  \lead{Reference} the environment of \code{erase\_global\_deps}
  (\code{erasure/theories/ErasureFunction.v}).
\end{lemma}
\begin{proof}
  \uses{def:Erasure-collectDeps}
  \leanok
  \lead{Idea} \code{hcollect} by kernel evaluation (\code{rfl}); the route by unfolding
  \code{eraseEntry} and \code{route} at \code{hcollect}.
\end{proof}

\begin{lemma}[Unfolding the traversal at the pure backend]
  \label{lem:EraseProof-Test-NV1-oracle_keep}
  \lean{EraseProof.Test.NV1.oracle_keep, EraseProof.Test.NV1.visitExpr_fvar_run,
    EraseProof.Test.NV1.visitExpr_lam_run, EraseProof.Test.NV1.visitAppArgs_one_run,
    EraseProof.Test.NV1.visitExpr_appFVar_run, EraseProof.Test.NV1.visitConst_run,
    EraseProof.Test.NV1.visitExpr_constApp_run, EraseProof.Test.NV1.get_constant_kername_run,
    EraseProof.Test.NV1.visitMutual_one_run}
  \leanok
  \uses{def:Erasure-erase, def:Erasure-PureM}
  \srcloc{proof/EraseProof/Test/NV1Run.lean}{46}
  \lead{In short} One equation per traversal call: from the runs of its callees and \hl{the
  oracle's answer \code{false}}, the call returns the erased term.
  \begin{itemize}
  \item \code{oracle\_keep}: after the answer \code{false}, \code{visitExpr} continues with its
        case branch;
  \item \code{visitExpr\_\{fvar, lam, appFVar, constApp\}\_run}, \code{visitAppArgs\_one\_run},
        \code{visitConst\_run}: a variable, a $\lambda$, an application of a variable or of a
        constant;
  \item \code{get\_constant\_kername\_run}: an unregistered constant runs \code{visitMutual} and
        returns its kername;
  \item \code{visitMutual\_one\_run}: \code{one}'s value is erased and registered as a constant.
  \end{itemize}
  \lead{Reference} the \code{tVar}, \code{tLambda}, \code{tApp}, \code{tConst} cases of
  \code{erase}, and \code{erase\_global\_deps} (\code{erasure/theories/ErasureFunction.v}).
\end{lemma}
\begin{proof}
  \uses{lem:EraseProof-visitLambda_agree, lem:EraseProof-visitExpr_box,
    lem:EraseProof-visitExpr_lam_step, lem:EraseProof-visitMutual_nonrec_step}
  \leanok
  \lead{Idea} Rewrite with the traversal's equation lemmas at \code{PureM} and the run lemmas of
  the frame and step nodes (\code{runPure\_bind}, \code{isErasable\_run}, \code{fresh\_run},
  \code{mkLambda\_run}, \code{name\_occurs\_eq}); the state maps are read by
  \code{getElem?\_insert\_self}.
\end{proof}

\begin{lemma}[NV-1: the oracle's answers]
  \label{lem:EraseProof-Test-NV1-oracle_oneE}
  \lean{EraseProof.Test.NV1.oracle_oneE, EraseProof.Test.NV1.oracle_lamZ,
    EraseProof.Test.NV1.oracle_app, EraseProof.Test.NV1.oracle_s, EraseProof.Test.NV1.oracle_z,
    EraseProof.Test.NV1.oracle_e0, EraseProof.Test.NV1.oracle_lamA,
    EraseProof.Test.NV1.oracle_a}
  \leanok
  \uses{def:Erasure-Pure-isErasable, def:EraseProof-Test-NV1-decls0, def:EraseProof-Test-NV1-p0}
  \srcloc{proof/EraseProof/Test/NV1Run.lean}{233}
  \lead{In short} The pure oracle \hl{keeps every term the run visits}: \code{one}'s value and
  its subterms, \code{e0}, and the argument \code{fun a => a} with its body.
\end{lemma}
\begin{proof}
  \uses{def:Erasure-Pure-isErasable}
  \leanok
  \lead{Idea} Kernel evaluation of the oracle (\code{rfl}).
\end{proof}

\begin{lemma}[NV-1: the run of \texorpdfstring{\code{\#erase}}{\#erase}]
  \label{lem:EraseProof-Test-NV1-hrun}
  \lean{EraseProof.Test.NV1.hrun, EraseProof.Test.NV1.hpure, EraseProof.Test.NV1.e0_run,
    EraseProof.Test.NV1.kername_one_run, EraseProof.Test.NV1.oneE_run,
    EraseProof.Test.NV1.body_lamZ_run, EraseProof.Test.NV1.body_app_run,
    EraseProof.Test.NV1.arg_run, EraseProof.Test.NV1.abstract_oneL}
  \leanok
  \uses{def:Erasure-eraseEntry, def:Erasure-erasePure, def:EraseProof-Test-NV1-decls0,
    def:EraseProof-Test-NV1-lenv0, def:EraseProof-Test-NV1-p0}
  \srcloc{proof/EraseProof/Test/NV1Run.lean}{329}
  \lead{In short} \code{hrun : eraseEntry view0 \{\} e0 = pure (p0, [])}: \hl{the hypothesis
  \code{hrun} of \code{erase\_correct}} holds on NV-1.
  \begin{itemize}
  \item \code{hpure}: \code{erasePure} over \code{decls0'} returns \code{(p0, [])};
  \item \code{e0\_run}: \code{visitExpr} on \code{e0} from the empty states returns \code{t0} and
        \code{stRun}, at every fuel from 14 on;
  \item the other names: the sub-runs, from \code{s z} up to \code{one}'s value and the argument.
  \end{itemize}
  \lead{Reference} \code{erase} and \code{erase\_global\_deps}
  (\code{erasure/theories/ErasureFunction.v}).
\end{lemma}
\begin{proof}
  \uses{lem:EraseProof-Test-NV1-hcollect, lem:EraseProof-Test-NV1-oracle_keep,
    lem:EraseProof-Test-NV1-oracle_oneE}
  \leanok
  \lead{Idea} Compose the unfolding lemmas bottom-up at the oracle's answers; \code{travFuel} is
  \code{(travFuel - 14) + 14}.
\end{proof}

\begin{lemma}[NV-1: the instance of \texorpdfstring{\code{erase\_correct}}{erase\_correct}]
  \label{lem:EraseProof-Test-NV1-final}
  \lean{EraseProof.Test.NV1.final}
  \leanok
  \uses{thm:EraseProof-erase_correct, def:EraseProof-Test-NV1-p0, def:EraseProof-Test-NV1-decls0,
    def:EraseProof-Test-NV1-lenv0, def:EraseProof-Erases, def:EraseProof-LBEval}
  \srcloc{proof/EraseProof/Test/NV1Run.lean}{337}
  \inherited{L1, L2, L3, L4, L5, L6}
  \lead{In short} \code{final} applies \code{erase\_correct} to \code{henv}, \code{hview},
  \code{he}, \code{hrun}, \code{hev}, \hl{every hypothesis a checked term}.
  \begin{itemize}
  \item \lead{Then} \code{p0 = untyped lenv0 (some t0)}, \code{v0'} erases \code{v0}, and
        \code{t0} evaluates to \code{v0' = fun z => (fun a => a) z}, the only witness.
  \end{itemize}
\end{lemma}
\begin{proof}
  \uses{thm:EraseProof-erase_correct, lem:EraseProof-Test-NV1-henv,
    lem:EraseProof-Test-NV1-hview, lem:EraseProof-Test-NV1-hrun, lem:EraseProof-Test-NV1-concl}
  \leanok
  \lead{Idea} Apply the theorem, inject \code{p0}, and fix the value by \code{witness}; the
  footprint is the theorem's.
\end{proof}

\section{Non-vacuity instance NV-2}\label{sec:tests-nv2}

\code{erases\_correct} on \code{dbl two}, whose evaluation takes $\delta$, $\beta$ and $\zeta$.

\begin{definition}[The program of NV-2]
  \label{def:EraseProof-Test-NV2-decls2}
  \lean{EraseProof.Test.NV2.decls2, EraseProof.Test.NV2.e2, EraseProof.Test.NV2.tyE,
    EraseProof.Test.NV2.cnE, EraseProof.Test.NV2.endoE, EraseProof.Test.NV2.cnatE,
    EraseProof.Test.NV2.twoE, EraseProof.Test.NV2.spE, EraseProof.Test.NV2.dblBodyE,
    EraseProof.Test.NV2.dblE, EraseProof.Test.NV2.CNat_val, EraseProof.Test.NV2.two_val,
    EraseProof.Test.NV2.dbl_val, EraseProof.Test.NV2.ty1, EraseProof.Test.NV2.ty2,
    EraseProof.Test.NV2.vCN, EraseProof.Test.NV2.endoV, EraseProof.Test.NV2.zV,
    EraseProof.Test.NV2.cnatV, EraseProof.Test.NV2.ssz, EraseProof.Test.NV2.twoV,
    EraseProof.Test.NV2.spV, EraseProof.Test.NV2.dblV, EraseProof.Test.NV2.CNdef,
    EraseProof.Test.NV2.twodef, EraseProof.Test.NV2.dbldef, EraseProof.Test.NV2.envCN0,
    EraseProof.Test.NV2.envCN, EraseProof.Test.NV2.envTwo0, EraseProof.Test.NV2.envTwo,
    EraseProof.Test.NV2.envDbl0, EraseProof.Test.NV2.venv2, EraseProof.Test.NV2.e2',
    EraseProof.Test.NV2.ctxD, EraseProof.Test.NV2.view2, EraseProof.Test.NV2.σ2,
    EraseProof.Test.NV2.v2}
  \leanok
  \uses{def:EraseProof-evalEnvOf}
  \srcloc{proof/EraseProof/Test/NV2.lean}{75}
  \lead{In short} \code{def CNat : Type 1 := forall a : Type, (a -> a) -> a -> a}, \code{def two},
  \code{def dbl (n : CNat) : CNat := let m := n; fun a s z => m a s (m a s z)}, and the term
  \hl{\code{dbl two}}.
  \begin{itemize}
  \item \code{decls2}, \code{e2}: the declarations, newest first, and the term;
  \item \code{venv2}, \code{e2'}: their lean4lean model and translation, built step by step;
  \item \code{view2}, \code{sigma2}: the evaluation environment; \code{v2}: the value of \code{e2}.
  \end{itemize}
\end{definition}

\begin{definition}[The erased program of NV-2]
  \label{def:EraseProof-Test-NV2-lenv2}
  \lean{EraseProof.Test.NV2.lenv2, EraseProof.Test.NV2.twoL, EraseProof.Test.NV2.spL,
    EraseProof.Test.NV2.dblBodyL, EraseProof.Test.NV2.dblL, EraseProof.Test.NV2.t2,
    EraseProof.Test.NV2.v2'}
  \leanok
  \uses{def:GlobalDecl, def:LBTerm}
  \srcloc{proof/EraseProof/Test/NV2.lean}{376}
  \lead{In short} The $\lambda_\Box$ environment \code{lenv2} holds \code{two} and \code{dbl} with
  their erased bodies; \hl{the type argument of \code{m} erases to $\Box$}.
  \begin{itemize}
  \item \code{t2}: \code{dbl two}; \code{v2'}: \code{fun a s z => two' box s (two' box s z)}, with
        \code{two'} the erased body of \code{two}.
  \end{itemize}
\end{definition}

\begin{lemma}[NV-2: the model-side hypotheses]
  \label{lem:EraseProof-Test-NV2-henv}
  \lean{EraseProof.Test.NV2.henv, EraseProof.Test.NV2.he, EraseProof.Test.NV2.endoV_ty,
    EraseProof.Test.NV2.zV_ty, EraseProof.Test.NV2.cnatV_ty, EraseProof.Test.NV2.cn_unfold,
    EraseProof.Test.NV2.vCN_ty, EraseProof.Test.NV2.sz_ty, EraseProof.Test.NV2.ssz_ty,
    EraseProof.Test.NV2.twoV_ty, EraseProof.Test.NV2.m_ty, EraseProof.Test.NV2.mα_ty,
    EraseProof.Test.NV2.mαs_ty, EraseProof.Test.NV2.spV_ty, EraseProof.Test.NV2.dblV_ty,
    EraseProof.Test.NV2.trTy, EraseProof.Test.NV2.trEndo, EraseProof.Test.NV2.trAlpha,
    EraseProof.Test.NV2.trCnat, EraseProof.Test.NV2.trTwo, EraseProof.Test.NV2.trSp,
    EraseProof.Test.NV2.trDbl, EraseProof.Test.NV2.pCN, EraseProof.Test.NV2.pTwo}
  \leanok
  \uses{def:EraseProof-Test-NV2-decls2, def:EraseProof-ProgEnv, def:EraseProof-TrS}
  \srcloc{proof/EraseProof/Test/NV2.lean}{297}
  \lead{In short} \code{henv : ProgEnv decls2 venv2} and \code{he : TrS venv2 [] [] e2 e2'}; the
  model types \code{m a} \hl{through \code{CNat}'s defining equation}.
\end{lemma}
\begin{proof}
  \uses{def:EraseProof-ProgEnv}
  \leanok
  \lead{Idea} One \code{ProgEnv} rule per declaration; typing derivations in lean4lean's
  \code{HasType}, with one $\delta$-conversion of \code{CNat} for \code{two} and two for
  \code{dbl}.
\end{proof}

\begin{lemma}[NV-2: the relation-level instance]
  \label{lem:EraseProof-Test-NV2-concl}
  \lean{EraseProof.Test.NV2.concl, EraseProof.Test.NV2.inst, EraseProof.Test.NV2.witness,
    EraseProof.Test.NV2.hev,
    EraseProof.Test.NV2.hsub, EraseProof.Test.NV2.mem_names, EraseProof.Test.NV2.hinj,
    EraseProof.Test.NV2.lookup2, EraseProof.Test.NV2.hlc, EraseProof.Test.NV2.hblocks,
    EraseProof.Test.NV2.erTwo, EraseProof.Test.NV2.erAlpha, EraseProof.Test.NV2.erSp,
    EraseProof.Test.NV2.erDbl, EraseProof.Test.NV2.her, EraseProof.Test.NV2.hdeps,
    EraseProof.Test.NV2.lbev, EraseProof.Test.NV2.herv}
  \leanok
  \uses{thm:EraseProof-erases_correct, lem:EraseProof-Test-NV2-henv,
    def:EraseProof-Test-NV2-lenv2}
  \srcloc{proof/EraseProof/Test/NV2.lean}{480}
  \inherited{L1, L2, L3, L4, L5, L6}
  \lead{In short} \code{inst} applies \code{erases\_correct} on NV-2 with \hl{every hypothesis a
  checked term}; \code{concl} is its conclusion at the $\lambda$ \code{v2'}.
  \begin{itemize}
  \item \code{erTwo}, \code{erAlpha}, \code{erSp}, \code{erDbl}: the erasures of the bodies, the
        type argument by \code{Erases.box};
  \item \code{witness}: every witness of the conclusion is \code{v2'}.
  \end{itemize}
\end{lemma}
\begin{proof}
  \uses{thm:EraseProof-erases_correct, lem:EraseProof-Test-LBEval-deterministic}
  \leanok
  \lead{Idea} As NV-1.
\end{proof}

\section{Non-vacuity instance NV-3}\label{sec:tests-nv3}

\code{erases\_correct} on \code{uf one}, with \code{uf} unsafe and recursive through its own name:
the source evaluation takes \code{fixAtom} and \code{fixApp}, $\lambda_\Box$ evaluation
\code{LBEval.fix}.

\begin{definition}[The program of NV-3]
  \label{def:EraseProof-Test-NV3-decls3}
  \lean{EraseProof.Test.NV3.decls3, EraseProof.Test.NV3.e3, EraseProof.Test.NV3.cnT,
    EraseProof.Test.NV3.cnArr, EraseProof.Test.NV3.ufBody, EraseProof.Test.NV3.ufE,
    EraseProof.Test.NV3.uf_val, EraseProof.Test.NV3.cnC, EraseProof.Test.NV3.cnArrV,
    EraseProof.Test.NV3.ufV, EraseProof.Test.NV3.ufdef, EraseProof.Test.NV3.envUf0,
    EraseProof.Test.NV3.venv3, EraseProof.Test.NV3.e3', EraseProof.Test.NV3.view3,
    EraseProof.Test.NV3.σ3}
  \leanok
  \uses{def:EraseProof-Test-NV1-decls0, def:EraseProof-evalEnvOf}
  \srcloc{proof/EraseProof/Test/NV3.lean}{58}
  \lead{In short} NV-1's program with \code{unsafe def uf (n : CN) : CN := (fun (k : CN -> CN) =>
  n) (fun (y : CN) => uf y)}, and the erased term \hl{\code{uf one}}.
  \begin{itemize}
  \item \code{decls3}, \code{e3}: the declarations, newest first, and the term;
  \item \code{venv3}, \code{e3'}: NV-1's model with the one-member block \code{[uf]}, and the
        translation;
  \item \code{view3}, \code{sigma3}: a list view and its evaluation environment.
  \end{itemize}
\end{definition}

\begin{definition}[The erased program of NV-3]
  \label{def:EraseProof-Test-NV3-lenv3}
  \lean{EraseProof.Test.NV3.lenv3, EraseProof.Test.NV3.ufL, EraseProof.Test.NV3.defs3,
    EraseProof.Test.NV3.t3, EraseProof.Test.NV3.ufL'}
  \leanok
  \uses{def:GlobalDecl, def:LBTerm}
  \srcloc{proof/EraseProof/Test/NV3.lean}{265}
  \lead{In short} The $\lambda_\Box$ environment \code{lenv3} holds \code{one} and \hl{\code{uf} as
  the fixpoint \code{tFix [uf] 0}}, as \code{\#erase} emits it.
  \begin{itemize}
  \item \code{ufL}: the fixpoint's body, whose recursive call is the fixpoint variable;
        \code{ufL'}: the body \code{cunfold\_fix} gives;
  \item \code{t3}: \code{uf one}, the erasure of \code{e3}.
  \end{itemize}
\end{definition}

\begin{lemma}[NV-3: the model-side hypotheses]
  \label{lem:EraseProof-Test-NV3-henv}
  \lean{EraseProof.Test.NV3.henv, EraseProof.Test.NV3.he, EraseProof.Test.NV3.cnC_ty,
    EraseProof.Test.NV3.cnArr_ty, EraseProof.Test.NV3.uf_ty, EraseProof.Test.NV3.one_ty,
    EraseProof.Test.NV3.envUf0_CN, EraseProof.Test.NV3.envUf0_uf, EraseProof.Test.NV3.venv3_CN,
    EraseProof.Test.NV3.venv3_uf, EraseProof.Test.NV3.venv3_one, EraseProof.Test.NV3.venv3_A,
    EraseProof.Test.NV3.kfun_ty, EraseProof.Test.NV3.yfun_ty, EraseProof.Test.NV3.ufV_ty,
    EraseProof.Test.NV3.trCNc, EraseProof.Test.NV3.trCnArr, EraseProof.Test.NV3.trUfE,
    EraseProof.Test.NV3.addUf}
  \leanok
  \uses{def:EraseProof-Test-NV3-decls3, def:EraseProof-ProgEnv, def:EraseProof-TrS}
  \srcloc{proof/EraseProof/Test/NV3.lean}{171}
  \lead{In short} \code{henv : ProgEnv decls3 venv3} and \code{he : TrS venv3 [] [] e3 e3'}; the
  value of \code{uf} is typed \hl{in the model that already declares \code{uf}}.
\end{lemma}
\begin{proof}
  \uses{lem:EraseProof-Test-NV1-henv}
  \leanok
  \lead{Idea} NV-1's \code{ProgEnv} extended by \code{ProgEnv.block} with one member; typing
  derivations in lean4lean's \code{HasType}.
\end{proof}

\begin{lemma}[NV-3: the relation-level instance]
  \label{lem:EraseProof-Test-NV3-inst}
  \lean{EraseProof.Test.NV3.inst, EraseProof.Test.NV3.concl, EraseProof.Test.NV3.witness,
    EraseProof.Test.NV3.hev, EraseProof.Test.NV3.hsub, EraseProof.Test.NV3.mem_names,
    EraseProof.Test.NV3.hinj, EraseProof.Test.NV3.atom_uf, EraseProof.Test.NV3.atom_one,
    EraseProof.Test.NV3.rec_uf, EraseProof.Test.NV3.rec_one, EraseProof.Test.NV3.unfold_uf,
    EraseProof.Test.NV3.unfold_one, EraseProof.Test.NV3.find_uf, EraseProof.Test.NV3.find_one,
    EraseProof.Test.NV3.substl_ufL, EraseProof.Test.NV3.cunfold, EraseProof.Test.NV3.look_uf,
    EraseProof.Test.NV3.look_one, EraseProof.Test.NV3.look_CN, EraseProof.Test.NV3.look_A,
    EraseProof.Test.NV3.recIn_uf, EraseProof.Test.NV3.herv, EraseProof.Test.NV3.erUf,
    EraseProof.Test.NV3.block_uf, EraseProof.Test.NV3.decl_uf, EraseProof.Test.NV3.decl_one,
    EraseProof.Test.NV3.deps_defs, EraseProof.Test.NV3.her, EraseProof.Test.NV3.hdeps,
    EraseProof.Test.NV3.hblocks, EraseProof.Test.NV3.look_cases, EraseProof.Test.NV3.hlc,
    EraseProof.Test.NV3.lbev}
  \leanok
  \uses{thm:EraseProof-erases_correct, lem:EraseProof-Test-NV3-henv,
    def:EraseProof-Test-NV3-lenv3}
  \srcloc{proof/EraseProof/Test/NV3.lean}{395}
  \inherited{L1, L2, L3, L4, L5, L6}
  \lead{In short} \code{inst} applies \code{erases\_correct} on NV-3 with every hypothesis a
  checked term; \hl{the rules \code{fixAtom} and \code{fixApp} have an instance}.
  \begin{itemize}
  \item \code{hev}: \code{fixAtom} evaluates the head \code{uf}, \code{fixApp} unfolds it once at
        \code{one}, two $\beta$ steps follow;
  \item \code{hblocks}, \code{block\_uf}: the one-member \code{ErasesBlock}; the recursive call
        erases by \code{Erases.constRec};
  \item \code{concl}, \code{witness}: the witness \code{oneL}, reached by \code{LBEval.fix}, is
        the only one.
  \end{itemize}
\end{lemma}
\begin{proof}
  \uses{thm:EraseProof-erases_correct, lem:EraseProof-Test-LBEval-deterministic,
    lem:EraseProof-Test-NV1-henv}
  \leanok
  \lead{Idea} As NV-1.
\end{proof}

\section{Non-vacuity instance NV-4}\label{sec:tests-nv4}

\code{erases\_correct} on \code{ua bfalse one.\{0\}}, with \code{ua}, \code{ub} a two-member unsafe
block that calls itself through the other member.

\begin{definition}[The program of NV-4]
  \label{def:EraseProof-Test-NV4-decls4}
  \lean{EraseProof.Test.NV4.decls4, EraseProof.Test.NV4.e4, EraseProof.Test.NV4.lu,
    EraseProof.Test.NV4.tyE, EraseProof.Test.NV4.endoE, EraseProof.Test.NV4.endo1E,
    EraseProof.Test.NV4.cnE, EraseProof.Test.NV4.cnC, EraseProof.Test.NV4.oneE,
    EraseProof.Test.NV4.spE, EraseProof.Test.NV4.succBodyE, EraseProof.Test.NV4.csuccE,
    EraseProof.Test.NV4.ty1E, EraseProof.Test.NV4.cbE, EraseProof.Test.NV4.cbC,
    EraseProof.Test.NV4.btrueE, EraseProof.Test.NV4.bfalseE, EraseProof.Test.NV4.cnArrE,
    EraseProof.Test.NV4.idE, EraseProof.Test.NV4.callE, EraseProof.Test.NV4.appSelE,
    EraseProof.Test.NV4.selE, EraseProof.Test.NV4.memberE, EraseProof.Test.NV4.memberTyE,
    EraseProof.Test.NV4.CN_val, EraseProof.Test.NV4.one_val, EraseProof.Test.NV4.csucc_val,
    EraseProof.Test.NV4.CB_val, EraseProof.Test.NV4.btrue_val, EraseProof.Test.NV4.bfalse_val,
    EraseProof.Test.NV4.ua_val, EraseProof.Test.NV4.ub_val, EraseProof.Test.NV4.vu,
    EraseProof.Test.NV4.lv2, EraseProof.Test.NV4.lv3, EraseProof.Test.NV4.vty,
    EraseProof.Test.NV4.endoV, EraseProof.Test.NV4.zV, EraseProof.Test.NV4.cnV,
    EraseProof.Test.NV4.vcn, EraseProof.Test.NV4.oneV, EraseProof.Test.NV4.spV,
    EraseProof.Test.NV4.csuccV, EraseProof.Test.NV4.ty2V, EraseProof.Test.NV4.cbV,
    EraseProof.Test.NV4.vcb, EraseProof.Test.NV4.btrueV, EraseProof.Test.NV4.bfalseV,
    EraseProof.Test.NV4.cn0, EraseProof.Test.NV4.cnArrV, EraseProof.Test.NV4.idV,
    EraseProof.Test.NV4.callV, EraseProof.Test.NV4.selV, EraseProof.Test.NV4.memberV,
    EraseProof.Test.NV4.memberTyV, EraseProof.Test.NV4.CNdef, EraseProof.Test.NV4.onedef,
    EraseProof.Test.NV4.csuccdef, EraseProof.Test.NV4.CBdef, EraseProof.Test.NV4.btruedef,
    EraseProof.Test.NV4.bfalsedef, EraseProof.Test.NV4.uadef, EraseProof.Test.NV4.ubdef,
    EraseProof.Test.NV4.withConst, EraseProof.Test.NV4.envCN, EraseProof.Test.NV4.envOne,
    EraseProof.Test.NV4.envCsucc, EraseProof.Test.NV4.envCB, EraseProof.Test.NV4.envBtrue,
    EraseProof.Test.NV4.envBfalse, EraseProof.Test.NV4.envBlk, EraseProof.Test.NV4.venv4,
    EraseProof.Test.NV4.e4', EraseProof.Test.NV4.view4, EraseProof.Test.NV4.σ4,
    EraseProof.Test.NV4.v4}
  \leanok
  \uses{def:EraseProof-evalEnvOf}
  \srcloc{proof/EraseProof/Test/NV4.lean}{166}
  \lead{In short} The corpus program \code{uaFalseOne}: universe-polymorphic Church numerals
  \code{CN}, \code{one}, \code{csucc}, Church booleans \code{CB}, \code{btrue}, \code{bfalse},
  and the mutual block \hl{\code{ua}, \code{ub}}.
  \begin{itemize}
  \item \code{decls4}, \code{e4}: the declarations, newest first, with the block as \code{ub},
        \code{ua} (the kernel's order), and the term;
  \item \code{venv4}, \code{e4'}: the model, built step by step, and the translation;
  \item \code{view4}, \code{sigma4}; \code{v4}: the value \code{fun a s z => s (one' a s z)}.
  \end{itemize}
\end{definition}

\begin{definition}[The erased program of NV-4]
  \label{def:EraseProof-Test-NV4-lenv4}
  \lean{EraseProof.Test.NV4.lenv4, EraseProof.Test.NV4.oneL, EraseProof.Test.NV4.spL,
    EraseProof.Test.NV4.succBodyL, EraseProof.Test.NV4.csuccL, EraseProof.Test.NV4.btrueL,
    EraseProof.Test.NV4.bfalseL, EraseProof.Test.NV4.idL, EraseProof.Test.NV4.callL,
    EraseProof.Test.NV4.appSelL, EraseProof.Test.NV4.memberL, EraseProof.Test.NV4.defs4,
    EraseProof.Test.NV4.t4, EraseProof.Test.NV4.v4'}
  \leanok
  \uses{def:GlobalDecl, def:LBTerm}
  \srcloc{proof/EraseProof/Test/NV4.lean}{868}
  \lead{In short} The $\lambda_\Box$ environment \code{lenv4} stores \code{ua} and \code{ub} as
  \hl{the members 0 and 1 of one fixpoint \code{fix defs4}}, with the erased bodies of the other
  constants.
  \begin{itemize}
  \item \code{t4}: \code{ua bfalse one}; \code{v4'}: \code{fun a s z => s (one' box s z)}, the
        erasure of \code{v4}.
  \end{itemize}
\end{definition}

\begin{lemma}[NV-4: the model-side hypotheses]
  \label{lem:EraseProof-Test-NV4-henv}
  \lean{EraseProof.Test.NV4.henv, EraseProof.Test.NV4.he, EraseProof.Test.NV4.vu_wf,
    EraseProof.Test.NV4.endoV_ty, EraseProof.Test.NV4.zV_ty, EraseProof.Test.NV4.cnV_level,
    EraseProof.Test.NV4.cnV_ty, EraseProof.Test.NV4.cn_unfold, EraseProof.Test.NV4.vcn_ty,
    EraseProof.Test.NV4.sz_ty, EraseProof.Test.NV4.oneV_ty, EraseProof.Test.NV4.n_ty,
    EraseProof.Test.NV4.nα_ty, EraseProof.Test.NV4.nαs_ty, EraseProof.Test.NV4.nαsz_ty,
    EraseProof.Test.NV4.csuccV_ty, EraseProof.Test.NV4.cbV_ty, EraseProof.Test.NV4.cb_unfold,
    EraseProof.Test.NV4.vcb_ty, EraseProof.Test.NV4.btrueV_ty, EraseProof.Test.NV4.bfalseV_ty,
    EraseProof.Test.NV4.cn0_ty, EraseProof.Test.NV4.cnArrV_ty, EraseProof.Test.NV4.idV_ty,
    EraseProof.Test.NV4.cbt_ty, EraseProof.Test.NV4.csucc0_ty, EraseProof.Test.NV4.csm_ty,
    EraseProof.Test.NV4.callBody_ty, EraseProof.Test.NV4.callV_ty, EraseProof.Test.NV4.b_ty,
    EraseProof.Test.NV4.bα_ty, EraseProof.Test.NV4.bαi_ty, EraseProof.Test.NV4.bαig_ty,
    EraseProof.Test.NV4.selV_ty, EraseProof.Test.NV4.memberV_ty, EraseProof.Test.NV4.memberTyV_ty,
    EraseProof.Test.NV4.trEndo, EraseProof.Test.NV4.trAlpha, EraseProof.Test.NV4.trCn,
    EraseProof.Test.NV4.trOne, EraseProof.Test.NV4.trCsucc, EraseProof.Test.NV4.trCb,
    EraseProof.Test.NV4.trBtrue, EraseProof.Test.NV4.trBfalse, EraseProof.Test.NV4.trCn0,
    EraseProof.Test.NV4.trCnArr, EraseProof.Test.NV4.trMemberTy, EraseProof.Test.NV4.trMember,
    EraseProof.Test.NV4.pCN, EraseProof.Test.NV4.pOne, EraseProof.Test.NV4.pCsucc,
    EraseProof.Test.NV4.pCB, EraseProof.Test.NV4.pBtrue, EraseProof.Test.NV4.pBfalse,
    EraseProof.Test.NV4.envBlk_dCB, EraseProof.Test.NV4.addBlk, EraseProof.Test.NV4.ua_bfalse_ty}
  \leanok
  \uses{def:EraseProof-Test-NV4-decls4, def:EraseProof-ProgEnv, def:EraseProof-TrS}
  \srcloc{proof/EraseProof/Test/NV4.lean}{677}
  \lead{In short} \code{henv : ProgEnv decls4 venv4} and \code{he : TrS venv4 [] [] e4 e4'};
  \code{ProgEnv.block} adds both members and types \hl{each value in a model where the other member
  is declared}.
\end{lemma}
\begin{proof}
  \uses{def:EraseProof-ProgEnv}
  \leanok
  \lead{Idea} One \code{ProgEnv} rule per declaration; typing derivations in lean4lean's
  \code{HasType}, with $\delta$-conversions of \code{CN} and \code{CB}.
\end{proof}

\begin{lemma}[NV-4: the source evaluation]
  \label{lem:EraseProof-Test-NV4-hev}
  \lean{EraseProof.Test.NV4.hev, EraseProof.Test.NV4.unfold_ua, EraseProof.Test.NV4.unfold_ub,
    EraseProof.Test.NV4.rec_ua, EraseProof.Test.NV4.rec_ub, EraseProof.Test.NV4.ev_bfalse,
    EraseProof.Test.NV4.ev_btrue, EraseProof.Test.NV4.ev_one0, EraseProof.Test.NV4.ev_ua_bfalse,
    EraseProof.Test.NV4.ev_ub_btrue, EraseProof.Test.NV4.ev_sel_false,
    EraseProof.Test.NV4.ev_sel_true, EraseProof.Test.NV4.ev_csucc, EraseProof.Test.NV4.ev_ub_call}
  \leanok
  \uses{def:EraseProof-Test-NV4-decls4, def:EraseProof-SrcEval}
  \srcloc{proof/EraseProof/Test/NV4.lean}{787}
  \lead{In short} \code{e4} evaluates to \code{v4}: \hl{\code{fixAtom} and \code{fixApp} on both
  members}, $\delta$ on \code{bfalse}, \code{btrue}, and at level 0 on \code{one} and
  \code{csucc}, and $\beta$.
\end{lemma}
\begin{proof}
  \uses{def:EraseProof-SrcEval}
  \leanok
  \lead{Idea} \code{bfalse} selects the branch that calls \code{ub} on \code{csucc one};
  \code{btrue} then selects the identity.
\end{proof}

\begin{lemma}[NV-4: the erasure hypotheses]
  \label{lem:EraseProof-Test-NV4-her}
  \lean{EraseProof.Test.NV4.her, EraseProof.Test.NV4.hdeps, EraseProof.Test.NV4.hblocks,
    EraseProof.Test.NV4.hlc, EraseProof.Test.NV4.mem_names, EraseProof.Test.NV4.cunfold_ua,
    EraseProof.Test.NV4.cunfold_ub, EraseProof.Test.NV4.look_one, EraseProof.Test.NV4.look_bfalse,
    EraseProof.Test.NV4.look_ub, EraseProof.Test.NV4.look_ua, EraseProof.Test.NV4.look_csucc,
    EraseProof.Test.NV4.look_btrue, EraseProof.Test.NV4.look_CB, EraseProof.Test.NV4.look_CN,
    EraseProof.Test.NV4.lookupConst_cons, EraseProof.Test.NV4.look_cases,
    EraseProof.Test.NV4.atom_ua, EraseProof.Test.NV4.atom_ub, EraseProof.Test.NV4.atom_bfalse,
    EraseProof.Test.NV4.atom_btrue, EraseProof.Test.NV4.atom_one, EraseProof.Test.NV4.atom_csucc,
    EraseProof.Test.NV4.recIn_ua, EraseProof.Test.NV4.recIn_ub, EraseProof.Test.NV4.erOne,
    EraseProof.Test.NV4.erAlpha, EraseProof.Test.NV4.erSuccBody, EraseProof.Test.NV4.erCsucc,
    EraseProof.Test.NV4.erBtrue, EraseProof.Test.NV4.erBfalse, EraseProof.Test.NV4.erCnArr,
    EraseProof.Test.NV4.erMember, EraseProof.Test.NV4.block4, EraseProof.Test.NV4.decl_ua,
    EraseProof.Test.NV4.decl_ub, EraseProof.Test.NV4.deps_btrue, EraseProof.Test.NV4.deps_csucc,
    EraseProof.Test.NV4.deps_memberL, EraseProof.Test.NV4.deps_defs, EraseProof.Test.NV4.deps_ua}
  \leanok
  \uses{def:EraseProof-Test-NV4-decls4, def:EraseProof-Test-NV4-lenv4, def:EraseProof-Erases,
    def:EraseProof-ErasesDeps, def:EraseProof-LenvClosed}
  \srcloc{proof/EraseProof/Test/NV4.lean}{1067}
  \lead{In short} \code{her}, \code{hdeps}, \code{hblocks} and \code{hlc} hold; \hl{the block
  satisfies the two-member \code{ErasesBlock}} (\code{block4}).
  \begin{itemize}
  \item \code{her}: \code{ua} erases to its kername, not to its fixpoint;
  \item \code{erMember}: a member's value erases to its fixpoint body, the call through
        \code{Erases.constRec}, the type argument to $\Box$.
  \end{itemize}
\end{lemma}
\begin{proof}
  \uses{def:EraseProof-Erases, lem:EraseProof-Test-NV4-henv, lem:EraseProof-Test-NV4-hev}
  \leanok
  \lead{Idea} Rule by rule; each stored body by its lookup lemma; \code{hlc} by
  \code{look\_cases}.
\end{proof}

\begin{lemma}[NV-4: the relation-level instance]
  \label{lem:EraseProof-Test-NV4-inst}
  \lean{EraseProof.Test.NV4.inst, EraseProof.Test.NV4.concl, EraseProof.Test.NV4.witness,
    EraseProof.Test.NV4.hsub, EraseProof.Test.NV4.hinj, EraseProof.Test.NV4.lb_ua_bfalse,
    EraseProof.Test.NV4.lb_ub_btrue, EraseProof.Test.NV4.lb_sel_false,
    EraseProof.Test.NV4.lb_sel_true, EraseProof.Test.NV4.lb_csucc, EraseProof.Test.NV4.lbev,
    EraseProof.Test.NV4.herv}
  \leanok
  \uses{thm:EraseProof-erases_correct, lem:EraseProof-Test-NV4-henv, lem:EraseProof-Test-NV4-hev,
    lem:EraseProof-Test-NV4-her}
  \srcloc{proof/EraseProof/Test/NV4.lean}{1152}
  \inherited{L1, L2, L3, L4, L5, L6}
  \lead{In short} \code{inst} applies \code{erases\_correct} on NV-4 with every hypothesis a
  checked term; $\lambda_\Box$ evaluation takes \hl{\code{LBEval.fix} on both members}.
  \begin{itemize}
  \item \code{hsub}, \code{hinj}: every declaration is the first of its name; the kernames are
        distinct;
  \item \code{concl}, \code{witness}: the witness \code{v4'} is the only one.
  \end{itemize}
\end{lemma}
\begin{proof}
  \uses{thm:EraseProof-erases_correct, lem:EraseProof-Test-LBEval-deterministic}
  \leanok
  \lead{Idea} As NV-1.
\end{proof}

\section{The statement of the final theorem}\label{sec:tests-stub}

The approved statement of \code{erase\_correct}, written out as a test, pins the theorem's type.

\begin{lemma}[The approved statement]
  \label{lem:EraseProof-Test-Stub-erase_correct}
  \lean{EraseProof.Test.Stub.erase_correct}
  \leanok
  \uses{def:EraseProof-ProgEnv, def:EraseProof-ViewAgrees, def:EraseProof-TrS,
    def:Erasure-eraseEntry, def:EraseProof-SrcEval, def:EraseProof-evalEnvOf,
    def:EraseProof-Erases, def:EraseProof-LBEval, def:EraseProof-defaultFlags}
  \srcloc{proof/EraseProof/Test/Stub.lean}{18}
  \inherited{L1, L2, L3, L4, L5, L6}
  \lead{In short} The statement of \code{erase\_correct} as approved, with the binders of its
  section, \hl{proved by \code{erase\_correct}}.
  \begin{itemize}
  \item \lead{Why} Check C12 of \code{proof/scripts/check.sh} fails unless this test and
        \code{erase\_correct} have the same type (\code{Expr.eqv}), binder order included.
  \end{itemize}
\end{lemma}
\begin{proof}
  \uses{thm:EraseProof-erase_correct}
  \leanok
  \lead{Idea} \code{erase\_correct} applied to the five hypotheses.
\end{proof}

\section{Bridge tests}\label{sec:tests-bridge}

\begin{lemma}[\texorpdfstring{\code{TrS}}{TrS} is an instance of \texorpdfstring{\code{TrExprS}}{TrExprS}]
  \label{lem:EraseProof-Test-TrS-toTrExprS}
  \lean{EraseProof.Test.TrS.toTrExprS}
  \leanok
  \uses{def:EraseProof-TrS}
  \srcloc{proof/EraseProof/Test/Bridge.lean}{23}
  \inherited{TrProj}
  \lead{In short} Every derivation of \code{TrS} is \hl{one of lean4lean's \code{TrExprS}}
  (DV-2, Section~\ref{reg:dv-2}).
\end{lemma}
\begin{proof}
  \uses{def:EraseProof-TrS}
  \leanok
  \lead{Idea} Induction, rule for rule.
\end{proof}

\begin{lemma}[\texorpdfstring{\code{ProgEnv}}{ProgEnv} is an instance of \texorpdfstring{\code{TrEnv'}}{TrEnv'}]
  \label{lem:EraseProof-Test-ProgEnv-toTrEnv}
  \lean{EraseProof.Test.ProgEnv.toTrEnv', EraseProof.Test.Bridge.ProgEnv.toTrEnv'_hashMap,
    EraseProof.Test.Bridge.TrConst.toTrConstant, EraseProof.Test.Bridge.TrDef.toTrDefVal,
    EraseProof.Test.Bridge.addConsts_fresh, EraseProof.Test.Bridge.constants_none_of_le,
    EraseProof.Test.Bridge.foldl_insert_of_forall_ne, EraseProof.Test.Bridge.foldl_insert_of_mem,
    EraseProof.Test.Bridge.fresh_of_addConst, EraseProof.Test.Bridge.insertDefs_hashMap,
    EraseProof.Test.Bridge.insert_step, EraseProof.Test.Bridge.unsafe_le}
  \leanok
  \uses{def:EraseProof-ProgEnv, lem:EraseProof-Test-TrS-toTrExprS}
  \srcloc{proof/EraseProof/Test/Bridge.lean}{196}
  \inherited{L1, TrProj}
  \lead{In short} A program environment is related by lean4lean's \code{TrEnv'} at safety
  \code{unsafe} to \hl{a constant map containing every declaration} (DV-3,
  Section~\ref{reg:dv-3}).
\end{lemma}
\begin{proof}
  \uses{lem:EraseProof-Test-TrS-toTrExprS, lem:EraseProof-ProgEnv-lookup}
  \leanok
  \lead{Idea} Induction on \code{ProgEnv}, building the map by insertions into a hash map.
\end{proof}

\section{Determinism of \texorpdfstring{$\lambda_\Box$}{lambda-box} evaluation}\label{sec:tests-lbeval}

\begin{lemma}[Determinism of \texorpdfstring{$\lambda_\Box$}{lambda-box} evaluation]
  \label{lem:EraseProof-Test-LBEval-deterministic}
  \lean{EraseProof.Test.LBEval.deterministic, EraseProof.Test.LBEval.spineArgs,
    EraseProof.Test.LBEval.headOf_mkApps, EraseProof.Test.LBEval.spineArgs_mkApps,
    EraseProof.Test.LBEval.mkApps_inj, EraseProof.Test.LBEval.mkApps_fix_inj,
    EraseProof.Test.LBEval.mkApps_construct_inj, EraseProof.Test.LBEval.forall₂_det,
    EraseProof.Test.LBEval.lbAtom_construct}
  \leanok
  \uses{def:EraseProof-LBEval}
  \srcloc{proof/EraseProof/Test/LBEval.lean}{90}
  \lead{In short} $\lambda_\Box$ evaluation is \hl{deterministic}, for every flag setting; the
  instances use it to show that the final theorem's witness is the one exhibited.
  \begin{itemize}
  \item \code{spineArgs}, \code{mkApps\_inj} and its variants: application spines are injective
        in their head and arguments;
  \item \code{forall}$_2$\code{\_det}: argument lists evaluated pointwise by a deterministic
        relation are equal (the \code{constructBlock} case);
  \item \code{lbAtom\_construct}: under \code{with\_constructor\_as\_block}, no constructor is an
        atom.
  \end{itemize}
  \lead{Reference} \code{eval\_deterministic} (\code{erasure/theories/EWcbvEval.v}).
\end{lemma}
\begin{proof}
  \uses{def:EraseProof-LBEval, lem:EraseProof-LBEval-ind}
  \leanok
  \lead{Idea} Induction on the first evaluation (\code{LBEval.ind}), inverting the second; the
  head tests of \code{appCong} separate it from the other application rules.
\end{proof}

\section{Register tests}\label{sec:tests-register}

The divergence register cites these tests. Each runs the shipping oracle at
\code{Erasure.oracleFuel} on a formal environment and is checked by \code{decide}.

\begin{lemma}[Register test: definition heads are not atoms]
  \label{lem:EraseProof-Test-Atoms-defHead_kept}
  \lean{EraseProof.Test.Atoms.defHead_kept, EraseProof.Test.Atoms.DefHead.ty0E,
    EraseProof.Test.Atoms.DefHead.ty1E, EraseProof.Test.Atoms.DefHead.AE,
    EraseProof.Test.Atoms.DefHead.A_val, EraseProof.Test.Atoms.DefHead.a_val,
    EraseProof.Test.Atoms.DefHead.Q_val, EraseProof.Test.Atoms.DefHead.hq_val,
    EraseProof.Test.Atoms.DefHead.decls, EraseProof.Test.Atoms.DefHead.view,
    EraseProof.Test.Atoms.DefHead.σ, EraseProof.Test.Atoms.DefHead.cx,
    EraseProof.Test.Atoms.DefHead.e0}
  \leanok
  \uses{def:Erasure-Pure-isErasable, def:EraseProof-EvalEnv-isAtom, def:EraseProof-evalEnvOf}
  \srcloc{proof/EraseProof/Test/Atoms.lean}{60}
  \lead{In short} With \code{def Q : Prop := forall P : Prop, P -> P} and \code{axiom hq : Q}, the
  oracle \hl{keeps the ill-typed spine \code{hq A a}}, and \code{hq} is not an atom: a proof
  headed by a definition cannot be one (DV-11).
\end{lemma}
\begin{proof}
  \uses{def:Erasure-Pure-isErasable}
  \leanok
  \lead{Idea} \code{decide}.
\end{proof}

\begin{lemma}[Register test: atomhood does not depend on levels]
  \label{lem:EraseProof-Test-Atoms-levelDependent_kept}
  \lean{EraseProof.Test.Atoms.levelDependent_kept, EraseProof.Test.Atoms.LevelDep.P_val,
    EraseProof.Test.Atoms.LevelDep.hq_val, EraseProof.Test.Atoms.LevelDep.decls,
    EraseProof.Test.Atoms.LevelDep.view, EraseProof.Test.Atoms.LevelDep.σ,
    EraseProof.Test.Atoms.LevelDep.cx}
  \leanok
  \uses{def:Erasure-Pure-isErasable, def:EraseProof-EvalEnv-isAtom, def:EraseProof-evalEnvOf}
  \srcloc{proof/EraseProof/Test/Atoms.lean}{90}
  \lead{In short} For \code{axiom hq.\{v\} : P.\{v\}} the oracle keeps \code{hq.\{v\}} and boxes
  \code{hq.\{0\}}, and \code{hq} is not an atom: \hl{atomhood cannot depend on the occurrence's
  levels} (DV-11).
\end{lemma}
\begin{proof}
  \uses{def:Erasure-Pure-isErasable}
  \leanok
  \lead{Idea} \code{decide}.
\end{proof}

\begin{lemma}[Register test: the oracle's \texorpdfstring{$\delta$}{delta}]
  \label{lem:EraseProof-Test-Oracle-irreducibleAlias}
  \lean{EraseProof.Test.Oracle.irreducibleAlias, EraseProof.Test.IrrAlias.ty1E,
    EraseProof.Test.IrrAlias.AE, EraseProof.Test.IrrAlias.A_val, EraseProof.Test.IrrAlias.IProp_val,
    EraseProof.Test.IrrAlias.R_val, EraseProof.Test.IrrAlias.hR_val,
    EraseProof.Test.IrrAlias.Endo_val, EraseProof.Test.IrrAlias.fI_val,
    EraseProof.Test.IrrAlias.a_val, EraseProof.Test.IrrAlias.decls, EraseProof.Test.IrrAlias.cx,
    EraseProof.Test.IrrAlias.guardE, EraseProof.Test.IrrAlias.appE, EraseProof.Test.IrrAlias.RE,
    EraseProof.Test.IrrAlias.hRE}
  \leanok
  \uses{def:Erasure-Pure-isErasable}
  \srcloc{proof/EraseProof/Test/Oracle.lean}{68}
  \lead{In short} Through the aliases \code{IProp := Prop} and \code{Endo := A -> A},
  \code{@[irreducible]} in the corpus programs, the oracle types and keeps \code{fun (\_ : R) (x :
  A) => x} and \code{fI a}, and \hl{boxes \code{R} and \code{hR}}, as the kernel's $\delta$ does
  (DV-18).
\end{lemma}
\begin{proof}
  \uses{def:Erasure-Pure-isErasable}
  \leanok
  \lead{Idea} \code{decide}.
\end{proof}

\section{Soundness instances of the oracle}\label{sec:tests-oracle}

\begin{definition}[The environment D1]
  \label{def:EraseProof-Test-D1-decls}
  \lean{EraseProof.Test.D1.decls, EraseProof.Test.D1.A_val, EraseProof.Test.D1.P_val,
    EraseProof.Test.D1.hP_val, EraseProof.Test.D1.F_val, EraseProof.Test.D1.view,
    EraseProof.Test.D1.σ, EraseProof.Test.D1.cx, EraseProof.Test.D1.vA, EraseProof.Test.D1.vP,
    EraseProof.Test.D1.ty1, EraseProof.Test.D1.AC, EraseProof.Test.D1.PC, EraseProof.Test.D1.hPC,
    EraseProof.Test.D1.FC, EraseProof.Test.D1.env1, EraseProof.Test.D1.env2,
    EraseProof.Test.D1.env3, EraseProof.Test.D1.env4}
  \leanok
  \uses{def:EraseProof-evalEnvOf}
  \srcloc{proof/EraseProof/Test/Oracle.lean}{100}
  \lead{In short} The axioms \hl{\code{A : Type}, \code{P : Prop}, \code{hP : P}, \code{F : Type ->
  Type}}, newest first, with their lean4lean model \code{env4}, built step by step
  (\code{env1}, ..., \code{env4}).
\end{definition}

\begin{lemma}[Instances of oracle soundness]
  \label{lem:EraseProof-Test-Oracle-sound_A}
  \lean{EraseProof.Test.Oracle.sound_A, EraseProof.Test.Oracle.sound_hP,
    EraseProof.Test.Oracle.oracle_A, EraseProof.Test.Oracle.oracle_hP,
    EraseProof.Test.Oracle.ok_of_toOption, EraseProof.Test.D1.henv, EraseProof.Test.D1.hsub,
    EraseProof.Test.D1.atom_A, EraseProof.Test.D1.atom_hP, EraseProof.Test.D1.atom_F,
    EraseProof.Test.D1.atom_P, EraseProof.Test.D1.e4A, EraseProof.Test.D1.e4P,
    EraseProof.Test.D1.e4hP, EraseProof.Test.D1.hPty2, EraseProof.Test.D1.hAty,
    EraseProof.Test.D1.hPty, EraseProof.Test.D1.hhPty}
  \leanok
  \uses{lem:EraseProof-Pure-isErasable_sound, def:EraseProof-Test-D1-decls}
  \srcloc{proof/EraseProof/Test/Oracle.lean}{204}
  \inherited{L1, L2, L3, L4, L5}
  \lead{In short} \code{Pure.isErasable\_sound} applied on D1 with no locals: \hl{every hypothesis is
  a checked term}, and the conclusions are that \code{A} and \code{hP} are erasable.
  \begin{itemize}
  \item \code{D1.henv : ProgEnv decls env4}, \code{D1.hsub : SubEnv decls decls}, the translations
        of \code{A} and \code{hP}, and the oracle's answers \code{oracle\_A}, \code{oracle\_hP}
        (by \code{decide});
  \item \code{atom\_A}, ...: the four axioms are atoms.
  \end{itemize}
\end{lemma}
\begin{proof}
  \uses{lem:EraseProof-Pure-isErasable_sound}
  \leanok
  \lead{Idea} Apply the theorem; its footprint is the instances' footprint.
\end{proof}

\section{Non-vacuity instances on atoms}\label{sec:tests-atoms}

\code{erases\_correct} on three programs that evaluate only through the atom rule
\code{constAtom} (DV-11). Each evaluates by $\beta$ to \code{fun x => x}, whose erasure is the
only witness.

\begin{lemma}[D1: the hypotheses on the environment]
  \label{lem:EraseProof-Test-D1-hinj}
  \lean{EraseProof.Test.D1.hinj, EraseProof.Test.D1.mem_names, EraseProof.Test.D1.hlc,
    EraseProof.Test.D1.hblocks}
  \leanok
  \uses{def:EraseProof-Test-D1-decls, def:EraseProof-KernameInj, def:EraseProof-LenvClosed,
    def:EraseProof-ErasesDeps}
  \srcloc{proof/EraseProof/Test/NV5.lean}{43}
  \lead{In short} On D1 with the empty $\lambda_\Box$ environment: \hl{distinct kernames}
  (\code{hinj}), a closed environment (\code{hlc}) and no stored fixpoint (\code{hblocks}).
\end{lemma}
\begin{proof}
  \uses{def:EraseProof-Test-D1-decls}
  \leanok
  \lead{Idea} \code{mem\_names} lists D1's four names; \code{decide} on each pair.
\end{proof}

\begin{definition}[The program of NV-5]
  \label{def:EraseProof-Test-NV5-e}
  \lean{EraseProof.Test.NV5.e, EraseProof.Test.NV5.idE, EraseProof.Test.NV5.v,
    EraseProof.Test.NV5.idV, EraseProof.Test.NV5.e', EraseProof.Test.NV5.t,
    EraseProof.Test.NV5.t'}
  \leanok
  \uses{def:EraseProof-Test-D1-decls, def:LBTerm}
  \srcloc{proof/EraseProof/Test/NV5.lean}{73}
  \lead{In short} \hl{\code{(fun (a : Type) (x : a) => x) A}} over D1: a type former as an
  argument; its image \code{e'}, its erasure \code{t}, the value \code{v} and its erasure
  \code{t'}.
\end{definition}

\begin{lemma}[NV-5: a type former as an argument]
  \label{lem:EraseProof-Test-NV5-inst}
  \lean{EraseProof.Test.NV5.inst, EraseProof.Test.NV5.concl, EraseProof.Test.NV5.witness,
    EraseProof.Test.NV5.hev, EraseProof.Test.NV5.idV_ty, EraseProof.Test.NV5.he,
    EraseProof.Test.NV5.her, EraseProof.Test.NV5.hdeps}
  \leanok
  \uses{thm:EraseProof-erases_correct, def:EraseProof-Test-NV5-e,
    lem:EraseProof-Test-D1-hinj, lem:EraseProof-Test-Oracle-sound_A}
  \srcloc{proof/EraseProof/Test/NV5.lean}{125}
  \inherited{L1, L2, L3, L4, L5, L6}
  \lead{In short} \code{inst} applies \code{erases\_correct} with \hl{every hypothesis a checked
  term}; \code{A} evaluates by \code{constAtom} and erases to $\Box$.
  \begin{itemize}
  \item \code{concl}: the conclusion with the witness \code{fun x => x}; \code{witness}: it is the
        only one.
  \end{itemize}
\end{lemma}
\begin{proof}
  \uses{thm:EraseProof-erases_correct, lem:EraseProof-Test-LBEval-deterministic}
  \leanok
  \lead{Idea} As NV-1; \code{D1.henv}, \code{D1.hsub} come from the oracle instances.
\end{proof}

\begin{definition}[The program of NV-6]
  \label{def:EraseProof-Test-NV6-e}
  \lean{EraseProof.Test.NV6.e, EraseProof.Test.NV6.kE, EraseProof.Test.NV6.v,
    EraseProof.Test.NV6.kV, EraseProof.Test.NV6.e', EraseProof.Test.NV6.t,
    EraseProof.Test.NV6.t'}
  \leanok
  \uses{def:EraseProof-Test-D1-decls, def:LBTerm}
  \srcloc{proof/EraseProof/Test/NV6.lean}{28}
  \lead{In short} \hl{\code{(fun (h : P) (y : Type) => y) hP}} over D1: a proof as an argument;
  its image, its erasure, its value and the value's erasure.
\end{definition}

\begin{lemma}[NV-6: a proof as an argument]
  \label{lem:EraseProof-Test-NV6-inst}
  \lean{EraseProof.Test.NV6.inst, EraseProof.Test.NV6.concl, EraseProof.Test.NV6.witness,
    EraseProof.Test.NV6.hev, EraseProof.Test.NV6.kV_ty, EraseProof.Test.NV6.he,
    EraseProof.Test.NV6.her, EraseProof.Test.NV6.hdeps, EraseProof.Test.NV6.hP_not_const}
  \leanok
  \uses{thm:EraseProof-erases_correct, def:EraseProof-Test-NV6-e,
    lem:EraseProof-Test-D1-hinj, lem:EraseProof-Test-Oracle-sound_A}
  \srcloc{proof/EraseProof/Test/NV6.lean}{81}
  \inherited{L1, L2, L3, L4, L5, L6}
  \lead{In short} \code{inst} applies \code{erases\_correct} with every hypothesis a checked term;
  the proof \code{hP} evaluates by \code{constAtom}.
  \begin{itemize}
  \item \code{concl}, \code{witness}: the witness \code{fun y => y}, the only one;
  \item \code{hP\_not\_const}: the atom \code{hP} \hl{erases only to $\Box$} (DV-11).
  \end{itemize}
\end{lemma}
\begin{proof}
  \uses{thm:EraseProof-erases_correct, lem:EraseProof-Test-LBEval-deterministic}
  \leanok
  \lead{Idea} As NV-5.
\end{proof}

\begin{definition}[The program of NV-7]
  \label{def:EraseProof-Test-NV7-e}
  \lean{EraseProof.Test.NV7.e, EraseProof.Test.NV7.decls, EraseProof.Test.NV7.A_val,
    EraseProof.Test.NV7.P_val, EraseProof.Test.NV7.hp_val, EraseProof.Test.NV7.view,
    EraseProof.Test.NV7.σ, EraseProof.Test.NV7.cx, EraseProof.Test.NV7.vA,
    EraseProof.Test.NV7.ty1, EraseProof.Test.NV7.vP0, EraseProof.Test.NV7.AC,
    EraseProof.Test.NV7.PC, EraseProof.Test.NV7.hpC, EraseProof.Test.NV7.env1,
    EraseProof.Test.NV7.env2, EraseProof.Test.NV7.env3, EraseProof.Test.NV7.v,
    EraseProof.Test.NV7.fE, EraseProof.Test.NV7.fV, EraseProof.Test.NV7.e',
    EraseProof.Test.NV7.t, EraseProof.Test.NV7.t'}
  \leanok
  \uses{def:EraseProof-evalEnvOf, def:LBTerm}
  \srcloc{proof/EraseProof/Test/NV7.lean}{147}
  \lead{In short} \code{(fun (h : P.\{0\}) (x : A) => x) hp} over the axioms \code{A : Type},
  \hl{\code{P.\{v\} : Sort v}}, \code{hp : P.\{0\}}, with the model \code{env3}.
\end{definition}

\begin{lemma}[NV-7: the hypotheses on the environment]
  \label{lem:EraseProof-Test-NV7-henv}
  \lean{EraseProof.Test.NV7.henv, EraseProof.Test.NV7.hsub, EraseProof.Test.NV7.mem_names,
    EraseProof.Test.NV7.hinj, EraseProof.Test.NV7.hlc, EraseProof.Test.NV7.hblocks,
    EraseProof.Test.NV7.atom_hp, EraseProof.Test.NV7.oracle_hp, EraseProof.Test.NV7.e3A,
    EraseProof.Test.NV7.e3P, EraseProof.Test.NV7.e3hp, EraseProof.Test.NV7.hP0ty,
    EraseProof.Test.NV7.hAty, EraseProof.Test.NV7.hhpty}
  \leanok
  \uses{def:EraseProof-Test-NV7-e, def:EraseProof-ProgEnv, def:EraseProof-EvalEnv-isAtom,
    def:Erasure-Pure-isErasable}
  \srcloc{proof/EraseProof/Test/NV7.lean}{96}
  \lead{In short} \code{henv}, \code{hsub}, \code{hinj}, \code{hlc}, \code{hblocks} hold, and
  \hl{\code{hp} is an atom} that the shipping oracle boxes, although \code{P}'s sort depends on a
  level parameter.
\end{lemma}
\begin{proof}
  \uses{def:EraseProof-ProgEnv}
  \leanok
  \lead{Idea} One \code{ProgEnv} rule per axiom; \code{atom\_hp} and \code{oracle\_hp} by
  \code{decide}.
\end{proof}

\begin{lemma}[NV-7: a proof at a universe-polymorphic proposition]
  \label{lem:EraseProof-Test-NV7-inst}
  \lean{EraseProof.Test.NV7.inst, EraseProof.Test.NV7.concl, EraseProof.Test.NV7.witness,
    EraseProof.Test.NV7.hev, EraseProof.Test.NV7.fV_ty, EraseProof.Test.NV7.he,
    EraseProof.Test.NV7.her, EraseProof.Test.NV7.hdeps}
  \leanok
  \uses{thm:EraseProof-erases_correct, lem:EraseProof-Test-NV7-henv}
  \srcloc{proof/EraseProof/Test/NV7.lean}{194}
  \inherited{L1, L2, L3, L4, L5, L6}
  \lead{In short} \code{inst} applies \code{erases\_correct} with every hypothesis a checked term;
  \hl{\code{hp} evaluates by \code{constAtom}} and erases to $\Box$.
  \begin{itemize}
  \item \code{concl}, \code{witness}: the witness \code{fun x => x}, the only one.
  \end{itemize}
\end{lemma}
\begin{proof}
  \uses{thm:EraseProof-erases_correct, lem:EraseProof-Test-LBEval-deterministic}
  \leanok
  \lead{Idea} As NV-5.
\end{proof}

\section{Non-vacuity instances}\label{sec:tests-instances}

Each instance discharges every hypothesis of a main statement by a checked term, in namespace
\code{EraseProof.Test.NV}\emph{k}: \code{inst} applies the statement, \code{concl} is its
conclusion at the exhibited witness, and \code{witness} shows that witness unique; NV-1's
instance of \code{erase\_correct} is \code{final}. The hypotheses, \code{concl} and \code{witness}
use the axioms \code{propext}, \code{Classical.choice}, \code{Quot.sound} only; \code{inst} and
\code{final} inherit their theorem's lean4lean sorries.

\begin{tabular}{lp{0.42\linewidth}p{0.36\linewidth}}
\toprule
Instance & Program & State \\
\midrule
NV-1 & \code{one (fun a : A => a)} & \stProved{} \code{erases\_correct}
  (Section~\ref{sec:tests-nv1}); \code{erase\_correct}, with \code{hrun} the run of
  \code{\#erase} (Section~\ref{sec:tests-nv1-run}) \\
NV-2 & \code{dbl two}, with a \code{let} in \code{dbl} & \stProved{} \code{erases\_correct}
  (Section~\ref{sec:tests-nv2}) \\
NV-3 & \code{uf one}, unsafe recursion & \stProved{} \code{erases\_correct}
  (Section~\ref{sec:tests-nv3}) \\
NV-4 & \code{ua bfalse one}, a two-member unsafe block & \stProved{} \code{erases\_correct}
  (Section~\ref{sec:tests-nv4}) \\
NV-5 & \code{(fun (a : Type) (x : a) => x) A} & \stProved{} \code{erases\_correct}
  (Section~\ref{sec:tests-atoms}) \\
NV-6 & \code{(fun (h : P) (y : Type) => y) hP} & \stProved{} \code{erases\_correct} \\
NV-7 & \code{(fun (\_ : P.\{0\}) (x : A) => x) hp} & \stProved{} \code{erases\_correct} \\
oracle & \code{A} and \code{hP} of D1 & \stProved{} both instances of
  \code{Pure.isErasable\_sound} (Section~\ref{sec:tests-oracle}) \\
\bottomrule
\end{tabular}

\section*{Caveats}
\begin{itemize}
\item \lead{Caveat} The bridge tests depend on \code{TrProj := sorry}; no other node may
  (Section~\ref{sec:trust-inherited}).
\item \lead{Caveat} Only NV-1 runs the eraser. NV-2 to NV-7 are relation-level: \code{her},
  \code{hdeps} and \code{hblocks} are exhibited by hand, not computed by \code{erasePure}.
\item \lead{Caveat} The kernel checks NV-1's run of \code{eraseEntry} over the list view
  \code{view0}. The step \code{nv1-emit} of \code{check.sh}, in CI, runs it over
  \code{EnvView.ofEnvironment} and compares the printing with \code{p0}'s; check C11 runs
  Peregrine on it, locally only.
\end{itemize}
